% Part: computability
% Chapter: computability-theory
% Section: def-functions-self-reference

\documentclass[../../../include/open-logic-section]{subfiles}

\begin{document}

\olfileid{cmp}{thy}{slf}
\olsection{స్వీయ-సూచనతో ప్రమేయాలను నిర్వచించడం}

ప్రమేయాలను వాటి పరంగానే నిర్వచించగలగడం సాధారణంగా ఉపయోగకరం.
ఉదాహరణకు, గణనీయ ప్రమేయాలు $k$, $l$,~$m$ ఇచ్చినప్పుడు, ప్రతి~$y$కు
కింది సమీకరణాన్ని నెరవేర్చే పాక్షిక గణనీయ ప్రమేయం~$f$ ఒకటి ఉందని
స్థిరబిందు ఉపపత్తి చెబుతుంది:
\[
f(y) \simeq
\begin{cases}
k(y) & \text{$l(y)=0$ అయితే} \\
f(m(y)) & \text{ఇతర సందర్భాల్లో.}
\end{cases}
\]
మరింత నిర్దిష్టంగా, $f$ను పొందడానికి
\[
g(x,y) \simeq
\begin{cases}
k(y) & \text{$l(y)=0$ అయితే} \\
\cfind{x}(m(y)) & \text{ఇతర సందర్భాల్లో}
\end{cases}
\]
అని ఉంచి, $\cfind{e}(y)=g(e,y)$ అయ్యే సూచిక~$e$ను కనుగొనడానికి
స్థిరబిందు ఉపపత్తిని వాడతాం.

ఒక నిర్దిష్ట ఉదాహరణగా, ``గరిష్ఠ సామాన్య భాజకం'' ప్రమేయం
$\fn{gcd}(u,v)$ను
\[
\fn{gcd}(u,v) \simeq
\begin{cases}
v & \text{$u=0$ అయితే} \\
\fn{gcd}(\fn{mod}(v,u),u) & \text{ఇతర సందర్భాల్లో}
\end{cases}
\]
అని నిర్వచించవచ్చు; ఇక్కడ $\fn{mod}(v,u)$ అనేది $v$ను~$u$తో
భాగించినప్పుడు వచ్చే శేషం. స్థిరబిందు ఉపపత్తిని ఆశ్రయిస్తే
$\fn{gcd}$ పాక్షిక గణనీయం అని వస్తుంది. (నిజానికి $y$ అనేది
$\tuple{u,v}$ జతను సంకేతీకరిస్తుందని తీసుకుని, దీన్ని పై రూపంలో
ఉంచవచ్చు.) ఆ తరువాత $u$పై ఆగమనం చేస్తే $\fn{gcd}$ నిజానికి
సర్వనిర్వచితమని తెలుస్తుంది.

ఇప్పుడే చర్చించిన ఉదాహరణల కంటే మరింత సంక్లిష్టమైన స్వీయ-సూచక
నిర్వచనాలను కూడా తప్పకుండా తయారు చేయవచ్చు. చాలా ప్రోగ్రామింగ్ భాషలు
ఏదో ఒక విధంగా ప్రమేయాలను వాటి పరంగానే నిర్వచించడాన్ని సమర్థిస్తాయి.
ఇది ఒక ప్రమేయాన్ని గణించే అల్గారిథానికి గల \emph{సూచిక} పరంగా ఆ
ప్రమేయాన్ని నిర్వచించగలగడం కంటే కొంచెం తక్కువ ఆశ్చర్యకరం అని గమనించండి;
స్థిరబిందు సిద్ధాంతం తన పూర్తి సాధారణతలో అనుమతించేది రెండవదే.

\end{document}
